\section{Matter and Radiation, One Tone}

There is \textbf{one tone}, not two. Matter and radiation are not two ingredients of the world. They are two
motions of the single ternary tone $\{-1, 0, +1\}$. Matter is the tone \emph{bound}. Radiation is the tone
\emph{waving}. There is no second field, no second alphabet, and no real numbers anywhere in the base.

\begin{principle}[Matter and radiation are one substance]
A standing knot is matter. A travelling ripple is radiation. They are the same tone-stuff in two motions. That
identity is mass-energy equivalence stated in the substrate.
\end{principle}

This section states the one-tone picture, the bind-versus-radiate problem it raised, and how soft radiation was
recovered with no real numbers. It closes with the capture-and-attraction program and its honest open edge. Two
terms recur, so a plain statement of each helps. A \emph{phonon} is a long-wavelength collective wave of the
medium, the kind of disturbance sound is made of. A \emph{bath} is an open or growing region a disturbance can
leave for good and never come back from. In this theory the bath is the wake, the ever-expanding frontier of the
mesh.

\subsection{The two motions of one tone}

There is one ternary tone field. Each site of a dock holds one tone in $\{-1, 0, +1\}$, and matter and radiation
are two kinds of excitation of that one field. They differ in their charge and their motion, not in their stuff.

\begin{table}[ht]
\centering
\begin{tabular}{@{}llll@{}}
\toprule
 & what it is & charge & motion \\
\midrule
matter & a localized bound charge (a particle) & net non-zero & trapped, circulating, can rest \\
radiation & a free propagating wave (a phonon) & net zero, $+$ and $-$ balance & always moving, massless \\
\bottomrule
\end{tabular}
\caption{Matter and radiation as two excitations of the one ternary tone field.}
\label{tab:two-motions}
\end{table}

Matter is the tone bound. Radiation is the tone waving. The bound body has rest mass because its energy is
trapped and confined. The wave is massless because a net-zero ripple has no rest center. Un-trap a bound
excitation and its mass turns into free radiation. Mass and radiation are one substance in two states, bound and
free.

\begin{principle}[Mass-energy equivalence in the substrate]
Rest mass is trapped energy. A bound circulating excitation has mass because its energy is confined. Release the
confinement and the mass becomes free radiation. This is $E = mc^2$ read off the one tone.
\end{principle}

\subsubsection{The drumhead}

The picture is a drumhead. Strike it once. It vibrates, radiates sound into the surrounding medium, and goes
still. The medium is not a second field. It is the same medium the drumhead sits in, and sound is just its waves.
Here the medium is the one tone field. Matter is a localized excitation of it, the drumhead. Radiation is a wave
of it, the sound. The bath is the open and growing mesh, and the sound radiates to the wake and never returns.
One field carries both. No second medium is needed.

\subsubsection{Why one field suffices}

Everything follows from the one tone and the one charge.

\begin{itemize}
\item \textbf{One charge.} Only the U(1) charge, the signed sum of tones in a dock, is conserved. Matter is a
region of net charge. A phonon is a net-zero ripple. There is no second charge.
\item \textbf{Neutral radiation.} A phonon carries equal plus and minus, so it carries a disturbance without
carrying net charge. The body keeps its identity, its charge, while shedding the disturbance.
\item \textbf{Settling.} Perturb the body's motion with a charge-neutral kick. It radiates that momentum away as
a phonon. The bath absorbs it. The body returns to its ground state.
\end{itemize}

That last point is the corrective agency of a self, achieved in one field.

\subsection{The wrong turn corrected}

An earlier line of reasoning concluded the substrate needed a \emph{second} ternary field, a $Z_3$ gauge field,
to carry the photon. That conclusion is wrong, and it violates the core commitment. The argument ran in three
steps.

\begin{enumerate}
\item Emission needs a neutral photon. (Correct.)
\item No fixed sector of the $24$ sites can be a static neutral photon, by the no-coloring enumeration.
(Correct.)
\item Therefore the photon must be a separate field. (Wrong.)
\end{enumerate}

Step three is the flaw. The statement that no fixed direction-sector is the photon only rules out the photon
being a \emph{static subset} of sites. It does not rule out the photon being a \emph{dynamical wave} of the one
tone.

\begin{result}[A wave is not a sector]
The no-coloring result rules out a photon that is a fixed subset of the $24$ sites. It says nothing about a
photon that is a propagating wave of the one tone. A wave is not a sector, so the enumeration does not touch it.
The second field is unnecessary, and one tone is kept.
\end{result}

The full $Z_3$ specification survives only as a richer-than-needed alternative recorded in an appendix. The
one-field photon-sector identity is owned by the gauge section.

\subsection{The bind-versus-radiate problem}

The one-tone picture asks for one thing. A single-field knit where the body \emph{confines}, so it has an
identity, but its \emph{disturbances propagate away} as phonons, so it can radiate and settle. This is the
\emph{leaky-breather} problem.

The naive way to confine a body is to \emph{reflect} its tones back inward at the body edge. That seals. The same
reflecting wall that bounces a lone charge inward also bounces the two halves of a neutral ripple, so the
disturbance the body must shed is trapped along with it. A reflecting confiner is dark to the bath. A self needs
both at once. A body genuinely bound, so a hit relaxes back, that still radiates the excess to the bath.

\begin{table}[ht]
\centering
\begin{tabular}{@{}lllp{0.30\textwidth}@{}}
\toprule
knit & confines a body & radiates to the bath & failure \\
\midrule
confinement by reflection & yes, extent $4$ & no, dark to the bath (cone $1$) & over-confines, seals its
radiation \\
free streaming & no, disperses (extent $9$) & yes, reaches the boundary (cone $7$) & no binding, a hit is never
healed \\
\bottomrule
\end{tabular}
\caption{The two simplest single-field knits each fail, for opposite reasons.}
\label{tab:bind-radiate}
\end{table}

Two attempted shortcuts were tested and failed honestly.

\begin{itemize}
\item \textbf{Co-motion is not binding.} A line of co-moving charges keeps formation only because lattice-gas
particles interact when they meet head-on, and parallel charges never meet. It is free parallel streaming, not
binding. The decisive test hits the body's own charges. Redirect one of the body's forward charges sideways and
ask whether the body pulls it back. On the closed torus, where nothing can leave so any healing would be real,
the redirected momentum flies off and never rejoins. Zero restoring force.
\item \textbf{The per-cell kink fails.} A topological kink, a domain wall in the tone, seemed bound by topology
rather than amplitude. But the same knit cannot give a clean kink and radiation at once. A radiating knit
shatters the kink, measured as one wall exploding into $135$. A static knit keeps the kink clean but traps a hit
in place forever.
\end{itemize}

\subsection{The root cause, then its undoing}

The kink failure pointed at a root cause. The ternary tone $\{-1, 0, +1\}$ looked too coarse for soft radiation.
The smallest per-site excitation is order one, so a disturbance is always order one. It either stays trapped or
becomes order-one turbulence. It never softly drains.

That reading was a red herring. It conflated two different things.

\begin{principle}[Softness is wavelength, not amplitude]
A soft mode is one of long wavelength, not small amplitude. The per-site amplitude can be coarse and order one
while the mode wavelength is unbounded. Softness lives at the collective scale, not the site scale.
\end{principle}

\begin{table}[ht]
\centering
\begin{tabular}{@{}lll@{}}
\toprule
quantity & ternary & matters for softness \\
\midrule
per-site amplitude & coarse, order one & no \\
mode wavelength & unbounded & yes \\
\bottomrule
\end{tabular}
\caption{Coarseness of the per-site value is irrelevant to softness. Wavelength is what matters.}
\label{tab:softness}
\end{table}

A soft mode is a low-frequency collective excitation spread over many docks, each still order one. Real matter
proves it. Atoms are discrete, yet a solid carries soft sound and heals deformations. A Boolean lattice gas, one
bit per link, recovers fluid dynamics, sound, and viscosity in the coarse-grained limit, the FHP result. One bit
is coarser than three tones and still enough.

\subsection{Soft radiation without real numbers}

The soft mode was built and measured on the committed substrate, with no real numbers anywhere. A coarse-grained
charge-density wave of wavelength $\lambda$ was built, deterministic and with zero net momentum. Its contrast
oscillates as the gas evolves. The half-period, the time to the first minimum, was read against the wavelength by
counting.

\begin{table}[ht]
\centering
\begin{tabular}{@{}lll@{}}
\toprule
wavelength $\lambda$ & half-period & speed $= \lambda / (2 \times \text{half-period})$ \\
\midrule
$2$ & $1$ & $1$ \\
$4$ & $2$ & $1$ \\
$6$ & $3$ & $1$ \\
$12$ & $6$ & $1$ \\
\bottomrule
\end{tabular}
\caption{The half-period grows linearly with wavelength at a constant speed of $1$.}
\label{tab:soft-dispersion}
\end{table}

The period grows linearly with wavelength at a constant speed of one.

\begin{result}[A gapless linear sound mode emerges from ternary tones]
The dispersion is $\omega = c\,k$, a gapless linear mode read entirely by counting. The frequency goes to zero as
the wavelength grows, which is exactly what soft means. The per-site coarseness is real but irrelevant. The soft
field lives at the collective scale.
\end{result}

So the radiation channel a self rides on is the emergent soft sound mode. The long-wavelength net part of a
disturbance leaves as sound. The short-wavelength part thermalizes.

\subsubsection{Two requirements that are not about discreteness}

Two further requirements hold for real-valued fields too, so they are not extra costs of going discrete.

\begin{itemize}
\item \textbf{Dissipation needs a bath or coarse-graining.} A reversible closed system cannot heal, it recurs, by
Poincar\'e. Healing needs energy to actually leave, either to the open and growing mesh, which is the bath, or to
fine degrees of freedom the coarse description drops.
\item \textbf{Clean hydrodynamics needs conservation laws and lattice symmetry.} Sound and viscosity emerge only
when the knit conserves the right quantities, charge and momentum, and the lattice is symmetric enough to be
isotropic. The committed model has a momentum-conserving collide step and the high-symmetry $24$ sites, so this
is met.
\end{itemize}

\subsection{Where the self lives}

The corrective self is an emergent structure, not a base per-site one. This is the committed picture. Vibes give
rise to emergent middle layers, and selves live there.

\begin{table}[ht]
\centering
\begin{tabular}{@{}llp{0.42\textwidth}@{}}
\toprule
aspect & level & basis \\
\midrule
identity & base & a topological charge, a kink's winding, conserved by the knit, no real numbers \\
agency & emergent & radiating a disturbance rides on the emergent soft sound mode \\
\bottomrule
\end{tabular}
\caption{Identity is hinted at the base, but corrective agency is emergent.}
\label{tab:where-self}
\end{table}

The reduced oscillator-bath model, a bound low-dimensional mode coupled to a soft radiative bath, is the correct
effective description of the emergent layer. The per-site test fails precisely because it probes below the scale
where the soft radiation can exist.

\subsection{The capture and attraction program}

Radiation is half of a self. The other half is \emph{attraction}, so bodies can bind and capture. The full ladder
was built, every step deterministic, on the committed sites where applicable, each with a control.

\begin{definition}[Capture]
A self is a captured composite seen at a coarse level. Capture is an attractive interaction together with a bath,
the attraction to draw bodies together and the bath to carry off the binding energy so they stay together.
\end{definition}

The single biggest lesson is that the bath enters twice. Once to enable capture, by shedding the binding energy.
Once to make the captured body into a self, by supplying the dissipation for an agential level.

\begin{table}[ht]
\centering
\begin{tabular}{@{}llp{0.40\textwidth}@{}}
\toprule
step & result & what it shows \\
\midrule
the bath exists & pass & the open mesh is a bath ($8 \to 0$), the torus is not ($8 \to 8$) \\
the recipe is correct & pass & capture needs attraction and a bath, both necessary \\
attraction emerges & pass & the lean's active vacuum mediates a Casimir attraction, no new ingredient \\
it captures & pass & mobile plates close $9 \to 2$ under the active vacuum, fixed under the inert one \\
reversible bound states & pass & a breather confines (extent $4$ vs $10$), reversible, but pinned \\
gauge binding & pass & the gauge force is constant, confining, strongest, but the coupling is an input \\
a captured body is not yet a self & pass (negative) & a reversible orbit has no degeneracy, no macro causation \\
the bath gives an attractor & pass & identity (spread $0.0001$) and agency (kick radiated, $0.0001$), a self \\
selves nest & pass & two bath-coupled bodies bind into a higher self, the construction recurses \\
\bottomrule
\end{tabular}
\caption{The capture and attraction ladder, each step deterministic with a control.}
\label{tab:capture-ladder}
\end{table}

The headline is that \textbf{attraction emerges from the base}, the eight atoms, via the lean's active vacuum. Two
walls suppress the vacuum density between them. The suppression falls with distance, $1.34$ at $d = 6$, $0.80$ at
$d = 8$, $0.23$ at $d = 10$, the Casimir hallmark. The wall-free control is exactly zero. No new ingredient, and
the bulk stays reversible.

This overturns an earlier reading. Two charges in the \emph{empty} vacuum genuinely do not interact. The
\emph{active} vacuum is a medium, and through it the knit mediates attraction. So the self program is not gated on
the gauge sector.

\subsection{The bind-versus-radiate resolution}

The tension is resolved, and the answer is one tone all the way down. There is no second field. The mistake was
confining the body by \emph{reflection}, which seals. The fix is to confine the body by a \emph{well}. There are
two ways to confine, and only one lets the body breathe.

\begin{table}[ht]
\centering
\begin{tabular}{@{}lp{0.24\textwidth}p{0.26\textwidth}l@{}}
\toprule
confine by & how it holds the body & what it does to a passing wave & result \\
\midrule
reflection & bounces lone charges inward at the body edge & bounces the halves of a neutral ripple at that same
edge & seals \\
a well & lowers the body's energy in a region, biasing its whole position & does not touch a passing wave, a
neutral ripple streams straight through & binds and radiates \\
\bottomrule
\end{tabular}
\caption{Two ways to confine. Only the well is transparent to radiation.}
\label{tab:well-vs-wall}
\end{table}

\begin{principle}[Bind by a well, not by a wall]
Bind the body's bulk by a well, not its individual tones by a wall. A well is transparent to radiation. A
reflecting wall is not. The body settles in the well while a neutral phonon passes straight through it and streams
to the bath.
\end{principle}

The reflecting confiner is the honest-negative contrast. It binds, but it is opaque to the very ripples the body
must shed, so it can never settle. The well binds the body's position while a neutral phonon streams to the bath.
Same one tone, bound by the well, waving as it leaves. \textbf{Matter binds and radiates from the single tone,
with no second alphabet.}

The well itself is not an added field. It is the inward radiation pressure of the depleted active vacuum. The
lean's create move fills empty space with zero-momentum tone pairs, a tone on a site and a tone on its opposite
site, which stream apart and later annihilate when they meet head-on. The body is mass, and mass excludes the
vacuum. The body absorbs the vacuum tones that stream into it, so it casts a shadow of reduced flux behind it.

\begin{theorem}[Shadow pressure binds, with the correct sign]
A body absorbs the active-vacuum tones streaming toward it, casting a shadow of reduced flux behind it. A
displaced piece is struck by more tones from the open side, which push it toward the body, than from the shadowed
body side, which push it away. The net push is toward the body. The binding force is attraction, fully discrete,
and reversible in the bulk.
\end{theorem}

The sign is the subtle point. The naive first-order reversible move, swapping the mass with the denser-side tone,
moves the mass toward higher density, which is repulsion, the wrong sign. The genuine force is not a swap. It is
the second-order pressure imbalance of the absorbed flux, and that imbalance is attraction. A crowd pushing on you
from one side pushes you the other way.

The deepest worry is that a body generating its own vacuum locally would refill its own shadow and erase the
force. It does not.

\begin{result}[The create move cannot wash out the shadow]
The create move makes zero-momentum pairs. A created pair carries momenta $\mathrm{root}(d)$ and
$\mathrm{root}(\overline{d}) = -\mathrm{root}(d)$, which cancel exactly. So the create move adds no net momentum
anywhere and cannot wash out a momentum imbalance. Only the body's asymmetric absorption sets a net momentum, so
even a maximally dense self-generated vacuum leaves the shadow's imbalance intact.
\end{result}

Measured on the real $24$-site dock, a self-contained body in its own dense active vacuum shows a net momentum at
a nearby plane of $-6$, pointing toward the body, the maximal shadow, against exactly zero with no body. The body
binds itself.

The mechanism is reversible because it is binding-by-radiation. The bulk streaming is an exact bijection, so it is
reversible. The body moves toward the depleted region by absorbing the asymmetric flux, and the energy it sheds
leaves as tones that stream to the wake and never return. The only irreversibility is absorption at the body and
at the bath, exactly where the model already quarantines it. This is why the open, growing mesh is non-negotiable.
A closed torus recurs, the shed energy comes back, and the body never settles. The bath makes
binding-by-radiation a one-way street.

\begin{principle}[Binding, radiating, and the bath are one act]
Binding is the inward radiation pressure of the depleted vacuum, the shadow. Radiation is the tones the body
absorbs and the disturbances it sheds. The bath is the wake those tones leave to and never return from, which
makes the settling one-way while the bulk stays reversible. Bind, radiate, and settle are the single act of
binding-by-radiation on the open mesh.
\end{principle}

\subsection{The honest open edge}

One gap remains, reported honestly rather than forced to a misleading pass. The committed experiments measure the
\emph{force}, the net momentum imbalance pointing toward the body, with a flat control. They show the well binds
and the reflecting wall seals. The remaining showcase is one long end-to-end run on the open mesh, a displaced
piece drifting home over many beats with a size sweep, and the same run recurring on the closed torus. The force
that would drive that movie is proven. The long-time settling movie is the remaining demonstration, not a
remaining mechanism. The mechanism is in hand, all on the one tone, fully discrete, reversible in the bulk.

\subsection{Why it matters}

Matter and radiation are not two ingredients. They are two motions of one tone. Matter is the tone bound.
Radiation is the tone waving. The bound body has rest mass because its energy is trapped. The wave is massless
because a net-zero ripple has no rest center. Soft radiation emerges as a collective long-wavelength mode, with no
real numbers. A self is a body that radiates its excess to the bath and returns to its ground state.

\begin{principle}[One field, one charge, bound or free]
One field. One charge. Bound or free. That is matter and radiation.
\end{principle}

In the felt image, a self is a knot of feeling in a sea of feeling. The sea is always stirring, pairs of pleasure
and pain welling up from peace and sinking back, everywhere and evenly. Where the knot sits the sea cannot also
be, so the knot stands in its own quiet, a hollow it carved by being. Because the sea presses in evenly from
everywhere except the quiet it made, it presses the knot's wandering pieces back home. The self is held together
by the very stillness it makes in the stirring world, and it stays itself by breathing the disturbance of being
out into the widening dark.
